% Part II: actions and jobs. Sources: docs/actions.md; docs/user-guide.md ("Actions and jobs");
% medorama/src/jobs.rs, actions.rs (secret_env), identify.rs; web/js/views/actions.js, jobs.js.
\chapter{Actions and jobs}
\label{chap:actions}

This chapter describes actions, the work that Medorama does on recordings in the background. It
is written for analysts who run actions, and for administrators and developers who install and
write plug-ins. An \emph{action} does something to one recording. Medorama's background workers
run it as a \emph{job} and keep its result. Two actions are built in, and \emph{plug-ins} add
others without a new build of Medorama. The \ui{Actions} tab runs actions on a selection of
recordings, and the \ui{Jobs} tab shows the queue.

\section{The built-in actions}
\label{sec:actions-builtin}

\begin{table}[htbp]
\caption{The built-in actions.}
\label{tab:actions-builtin}
\begin{tabularx}{\linewidth}{@{}P{3.0cm}L@{}}
\toprule
{\tabhead Action} & {\tabhead Does} \\
\midrule
\ui{Index} (\code{index}) & Builds the index that the viewers draw from (envelopes, markers,
  vital signs, events), in \file{index\_dir/<index key>/}, with the recording engine's WFDB or
  EDF reader, or with the recording's adapter (\cref{chap:formats-reference},
  \cref{chap:readers}). It is ticked by default in the Actions tab, and it runs automatically for
  new and changed recordings while the server runs (the setting \code{auto\_index}). \\
\ui{Read identifiers} (\code{identify}) & Reads the patient identifiers that the recording's own
  file states into the catalog's database only, so that they can be checked against the
  databases. Medorama reads the patient and recording fields of an EDF or BDF header itself;
  for other formats, the recording's adapter returns them. WFDB records hold none. The values
  are shown only to users with the PHI permission (\cref{chap:recordings}). This action also
  runs automatically. \\
\bottomrule
\end{tabularx}
\end{table}

\paragraph{A Holter's heart rate} is not an action of its own. The \code{index} job computes it
after indexing a recording that has an ECG lead and no heart rate of its own, such as a Holter
ECG or a sleep study. A Pan--Tompkins detector finds the QRS complexes, reading an hour of
signal at a time. For each minute, the job then keeps the median of the beat-to-beat rates and
their 5th to 95th percentile range; minutes with fewer than 10 plausible beats (30 to
250\,bpm) are left out. The series is kept in the database (\code{recording\_series}) and shown
on the patient page (\cref{chap:patients}).

\section{The Actions tab}
\label{sec:actions-tab}

The \ui{Actions} tab is shown to analysts and administrators. It lists the installed actions,
each with a box to tick, its title (built-in actions are marked \ui{built in}), its version, the
actions it requires, and its description. An action that cannot run (one that requires an action
that is not installed, or one of the reserved \code{encounter} scope) says why, and cannot be
ticked. Plug-ins whose manifest could not be read are listed after the table (``Not loaded:
\dots''). The ticked actions are remembered in this browser.

\ui{Run on} chooses the recordings:
\begin{description}
\item[\ui{The selected encounters}] The encounters ticked in the Patient Catalog's
  \ui{Encounters} table (\cref{chap:patients}), whose \ui{Run actions on selected\dots} button
  opens this tab.
\item[\ui{The encounters matching the current search}] The encounters that the Encounters
  table's search finds.
\item[\ui{The members of a cohort}] The recordings of a cohort's members, with the cohort chosen
  in the list. A cohort's \ui{Run actions\dots} opens this tab with that cohort chosen
  (\cref{chap:cohorts}).
\item[\ui{Every recording in the catalog}] All of them, after a confirmation.
\end{description}
\ui{Run again where already done} runs the ticked actions even where they already succeeded.
\ui{Run} queues the jobs and says what it did, for example: ``Run 12: 8 recording(s), 6 job(s)
queued; 2 already done, 0 already queued, 0 not applicable.''

\section{Runs and jobs}
\label{sec:actions-runs}

\paragraph{A run} applies one or more actions to a selection (encounters, a search, a cohort's
members, recordings, or all of them). It is expanded to one job per action and recording, with
the following rules:
\begin{itemize}
\item Missing recordings, and recordings of kinds that an action does not apply to, are left out
  (``not applicable'').
\item The actions that an action requires are added.
\item \textbf{No work is done twice.} A job is not queued when the same action, at the same
  version and with the same parameters, already succeeded on the recording as it is now (its
  fingerprint). \ui{Run again where already done} (\opt{-{}-force}) overrides this for the
  chosen actions.
\item A job already queued or running is never queued twice. Asking again with a higher
  priority raises its priority.
\end{itemize}

\paragraph{Workers} take the jobs from the queue. They are the workers of \cmd{medorama serve}
(as many at once as the setting \code{workers}, 2 by default), and \cmd{medorama work}, which
runs the queue in the terminal until it is empty. Workers take the job with the highest priority
first, then the oldest. \ui{Index now}, beside a recording, has priority 10; automatic indexing,
and runs unless asked otherwise, have priority 0. A job waits until the actions it requires have
succeeded on the recording, and fails if one of them failed.

\paragraph{Robustness.} A running job's worker writes a heartbeat every 2 seconds. A job whose
heartbeat stops for 2 minutes (a crashed or killed worker) goes back to the queue. A heartbeat
the database refuses (busy, a connection lost) is tried again: only the job's row stops a job
(cancelled, or taken back by the queue). A job's end is written again, for up to half a minute,
until the database takes it. A server or \cmd{medorama work} that stops (\key{Ctrl}+\key{C},
\code{SIGTERM}) stops its plug-ins, with everything they started (TERM to the plug-in's process
group, KILL 5 seconds later), and puts their jobs back in the queue at once. Canceling stops a
plug-in at once, while indexing runs to its end and its result is discarded. Several processes
(the server, \cmd{work}, the commands) can share one database, and each job is taken by exactly
one worker.

\section{The Jobs tab}
\label{sec:actions-jobs}

The \ui{Jobs} tab, shown to everyone, lists the jobs, unfinished ones first, or only those of one
state (\ui{Running}, \ui{Queued}, \ui{Failed}, \ui{Done}, \ui{Cancelled}), with the number of
jobs in each state. Each job shows its number, its state, its action, its recording (and its
error, when it failed), its progress while it runs, and when it was queued and finished. While
jobs are unfinished, the table refreshes itself every few seconds, and the number beside the
tab counts the jobs running and queued.

\begin{itemize}
\item \ui{Cancel}, beside a queued or running job, cancels it (analysts and administrators).
\item \ui{Cancel all unfinished} cancels every queued and running job, after a confirmation
  (administrators).
\end{itemize}

A failed job's error says why it failed. For a plug-in, it reads ``\dots{} exited with \dots''
followed by the last lines of its log, ``stopped after the time limit (\dots{} s)'', or
``canceled''. A recording whose indexing failed shows the error beside the recording too
(\cref{chap:patients}). The page does not show a job's result (its summary). The summary is kept
in the database with its recording, and the API returns it with the recording
(\code{GET recordings/\{id\}}) and with an encounter's recordings (\code{GET encounters/\{id\}};
\cref{chap:http-api}).

\section{From the command line}
\label{sec:actions-cli}

The following commands list the installed actions, queue runs, work through the queue, list
the failed jobs and cancel a job:
\begin{lstlisting}[style=shell]
medorama actions                               # the installed actions
medorama run index --all                       # queue
medorama run recording-summary --cohort 7 --param recording-summary.max_channels=16
medorama work                                  # run the queue here, until it is empty
medorama jobs --state failed                   # what failed, and why
medorama cancel --job 12                       # cancel a job (--run N: a run's, --all: every one)
\end{lstlisting}
\cmd{medorama run} takes the actions' names and a selection: \opt{-{}-all},
\opt{-{}-encounter ID} (repeated), \opt{-{}-recording ID} (repeated), \opt{-{}-search TEXT} (in
the encounters' search syntax; \cref{chap:patients}), or \opt{-{}-cohort N} (the cohort's
number, as the Cohorts tab shows it). \opt{-{}-force} runs the actions again where they are
already done, \opt{-{}-param ACTION.KEY=VALUE} overrides a parameter (VALUE is JSON, or text),
and \opt{-{}-priority N} sets the jobs' priority. The jobs are then run by \cmd{serve} or
\cmd{work}. \Cref{chap:cli} describes every command and option.

\section{Plug-ins}
\label{sec:actions-plugins}

A plug-in is a folder holding an \file{action.toml} file and a program, in one of the folders of
the setting \code{actions\_dirs} (by default \file{<home>/actions}). Medorama reads the plug-ins
when it starts, and again on \code{POST api/v1/actions/reload} (administrators). \cmd{serve}
prints the plug-ins it could not load, and the Actions tab lists the others. An example in
Python, using the standard library only, is in
\file{medorama/actions/examples/recording-summary}. Its \file{action.toml} file reads:
\begin{lstlisting}[language=TOML]
name = "recording-summary"
title = "Recording summary (example)"
version = "1"
description = "Channels, vital signs and duration of a recording, from its index: ..."
scope = "recording"
kinds = ["wfdb"]
requires = ["index"]
command = ["python3", "summary.py"]   # on Windows: ["python", "summary.py"]
timeout_s = 600
default = false

[params]
max_channels = 64                     # list at most this many channels in summary.json
\end{lstlisting}

\begin{xltabular}{\linewidth}{@{}P{2.6cm}P{2.6cm}L@{}}
\caption{The keys of \file{action.toml}.}
\label{tab:actions-manifest}\\
\toprule
{\tabhead Key} & {\tabhead Default} & {\tabhead Meaning} \\
\midrule
\endfirsthead
\toprule
{\tabhead Key} & {\tabhead Default} & {\tabhead Meaning} \\
\midrule
\endhead
\bottomrule
\endfoot
\code{name} & (Required) & The action's name: \code{a}--\code{z}, \code{0}--\code{9}, \code{\_}
  and \code{-}, up to 64 characters. It must be unique: a second action of the same name is not
  loaded. \\
\code{title} & The name & What the Actions tab shows. \\
\code{version} & \code{"1"} & A text or a number. Change it whenever the results would differ,
  so that the action runs again. \\
\code{description} & (Empty) & What the action does. \\
\code{scope} & \code{"recording"} & \code{"recording"}. The scope \code{"encounter"} is
  reserved: such actions are listed but cannot run yet. \\
\code{kinds} & Every kind & The recording kinds it applies to: \code{wfdb}, \code{edf}, or an
  adapter's kind (\code{["wfdb"]}). \\
\code{requires} & None & Actions that must have succeeded on the recording first. An action that
  requires one that is not installed cannot run. \\
\code{command} & (Required) & The program and its arguments. Relative paths that exist in the
  action's folder are made absolute. \\
\code{timeout\_s} & \code{21600} (6 hours) & The time limit, in seconds (0: no limit). \\
\code{default} & \code{false} & Whether the action is ticked in the Actions tab, until the user
  makes a choice. \\
\code{[params]} & None & Default parameters, which a run may override (\opt{-{}-param}). \\
\end{xltabular}

\paragraph{Trying the example.} Add the folder \file{medorama/actions/examples} of the clone
(its full path) to \code{actions\_dirs} in \file{medorama.toml} and restart the server. The
following command then queues the example on every recording and runs the queue:
\begin{lstlisting}[style=shell]
medorama run recording-summary --all && medorama work
\end{lstlisting}
On Windows, change its \code{command} to \code{["python", "summary.py"]}.

\section{The contract}
\label{sec:actions-contract}

For each job, Medorama:
\begin{enumerate}
\item creates the output folder
  \file{results\_dir/actions/<name>/r<recording id>/j<job id>/} (a job run again, after its
  worker stopped, gets a new one: \file{j<job id>-2/}, \file{-3/} \dots, the earlier attempt's
  files kept);
\item runs \code{command} there, with the job description as JSON on its standard input, also
  saved as \file{job.json} (its path in the environment variable \code{MEDORAMA\_JOB});
\item sends the command's standard error to \file{log.txt}, and keeps the other lines of its
  standard output in \file{stdout.txt}.
\end{enumerate}

The command must follow these rules:
\begin{itemize}
\item It writes its results into its working folder, never anywhere else.
\item It may print \code{progress <0..1>} lines on its standard output (\code{progress 0.5}),
  which the Jobs tab shows.
\item It may leave a \file{summary.json} file (up to 256\,KB), which Medorama stores in the
  database with the recording and returns through the API (\code{GET recordings/\{id\}}, and
  with an encounter's recordings through \code{GET encounters/\{id\}}). A larger summary is
  replaced by a note (``summary.json too large to store''), and one that is not JSON fails the
  job.
\item It exits with 0 on success. Any other exit, a timeout or a cancellation marks the job
  failed, and the failure message ends with the last lines of \file{log.txt}. At the time limit
  or a cancellation, the plug-in's process group gets TERM, then KILL 5 seconds later: a plug-in
  that must clean up handles TERM.
\end{itemize}

The command runs with the environment of the catalog's account, minus two groups of variables:
Medorama's own variables (\code{MEDORAMA\_\dots}: the database's address and password, the
encryption key), and every variable whose name marks it as secret (\code{PGPASS\dots}, names
containing \code{PASSWORD}, \code{SECRET}, \code{TOKEN}, \code{PASSPHRASE} or
\code{CREDENTIAL}, names ending in \code{\_KEY}, and \code{DATABASE\_URL}). The rest stays, so
that a plug-in finds its interpreter and libraries (\code{PATH}, \code{HOME}, a virtual or conda
environment). The command also receives \code{MEDORAMA\_JOB}, the path of its job description.

\paragraph{The job description} has this form:
\begin{lstlisting}[language=JSON]
{
  "catalog": "medorama 0.1.0",
  "job_id": 12, "action": "heart_rate", "version": "1",
  "params": {"window_s": 10},
  "recording": {"id": 3, "path": "/data/.../ICU7_12-...", "kind": "wfdb",
                "name": "ICU7_12-...", "format": "WFDB collection",
                "t_start": 1704106200.0, "t_end": 1704297600.0,
                "index_dir": "/results/index/0123456789abcdef0123"},
  "encounters": [{"encounter_id": "MO-E-7F3K-2QXM", "patient_id": "MO-AB4C-9D2E"}],
  "output_dir": "/results/actions/heart_rate/r3/j12"
}
\end{lstlisting}
\code{index\_dir} is \code{null} when the recording has no current index. An action that needs
one declares \code{requires = ["index"]} and reads it as the index format describes
(memory-mapped arrays and \file{meta.json}; \cref{chap:formats-reference}).

\section{Privacy}
\label{sec:actions-privacy}

\begin{phibox}
The job description holds pseudonyms only (\code{MO-\dots}), never identifiers, and job
summaries never contain identifiers (\code{identify} says only whether it found any). A plug-in
runs with the rights of the catalog's account. Install only plug-ins you trust, keep the actions
folders writable by that account only, and review what a plug-in writes to its results folder,
which is outside the catalog's database and its privacy gate (\cref{chap:security}).
\end{phibox}
